\scanpage{113}
% ms. line 01
Soit $(\alpha_n)$ une suite \uline{décroissante} \inserted{et convergeant vers zéro} de nombres
% ms. line 02
positifs. Soit $\rho$ la borne inférieure des nombres
% ms. line 03
\deleted{positifs} $s>0$ tels que $(\alpha_n)$ soit majorée par
% ms. line 04
un multiple de la série $\left(\frac{1}{n^{1/s}}\right)$ (et $+\infty$
% ms. line 05
si un tel $s$ n'existe pas) et soit $\sigma$ la
% ms. line 06
borne inférieure des nombres $s>0$, tels que
% ms. line 07
$\sum_n\alpha_n^s<+\infty$. (et $+\infty$ si un tel $s$ n'existe pas)

% ms. line 08
On a aussi $\frac1\rho\deleted{\unsure{\leq}}=\limsup_n\unsure{:}\frac{\log(1/\alpha_n)}{\log n}$. \deleted{$\rho$ et $\sigma$
% ms. line 09
sont resp\textsuperscript{t} les bornes inf. d'ensembles
% ms. line 10
A et B.} -- Le critère de Riemann affirme
% ms. line 11
que \deleted{$\unsure{\rho\leq|\rho\leq\sigma}$} $\sigma\leq\rho$. Mais on a aussi $\rho\leq\sigma$,
% ms. line 12
car on a le

% ms. line 13
\uline{Lemme} \quad Si $s>0$ est tel que $\sum_n\alpha_n^s<+\infty$
% ms. line 14
alors \inserted{$\unsure{\geq}$} \deleted{la suite $(\alpha_n)$ est majorée par un
% ms. line 15
multiple de la suite $\left(\frac1{n^{1/s}}\right)$.} \inserted{$\alpha_n=o\left(\frac1{n^{1/s}}\right)$} ($\alpha_n$ étant décroissant !)

% ms. line 16
Démonstration. On se ramène au cas où $s=1$
% ms. line 17
(en posant $B_n=\alpha_n^s$) \inserted{et à prouver qu'il est impossible que} \deleted{\unsure{sup}} pour une suite
% ms. line 18
strictement croissante d'entiers $n_i$, on \deleted{avait} ait
% ms. line 19
$\alpha_{n_i}\geq\frac1{n_i}$ ; \inserted{en effet, alors} on aurait \quad $\alpha_{n_{i-1}+1}+\alpha_{n_{i-1}+2}+\ldots+\alpha_{\unsure{n_i|n}}\geq\frac{n_i-n_{i-1}}{n_i}$
% ms. line 20
soit $\varepsilon_i$ ce second membre, $\sum\varepsilon_i$ est donc
% ms. line 21
convergente, et on a $n_i(1-\varepsilon_i)=n_{i-1}$, d'où $n_i=\frac{n_{i-1}}{1-\varepsilon_i}$
% ms. line 22
$\displaystyle =\frac{n_{i-2}}{(1-\varepsilon_i)(1-\varepsilon_{i-1})}\ldots=\frac1{(1-\varepsilon_{\unsure{1|i}})\ldots(1-\varepsilon_{\unsure{n|i}})}\leq\frac1{\prod_1^\infty(1-\varepsilon_k)}$, la suite
% ms. line 23
des $n_i$ est donc majorée, et s'arrête donc.
% archive sheet number: 114
